\documentclass[10pt]{beamer}
\usepackage[backend=biber,style=ieee]{biblatex}
\usepackage{tikz}
\usepackage[T1]{fontenc}
\usepackage{graphicx}
\usepackage{adjustbox}
\usepackage{minted}
\usepackage{booktabs}

\usepackage[frenchb]{babel}
\usepackage[babel]{csquotes}
\usepackage{adjustbox}
%\usepackage{lmodern}
%\usepackage{palatino}

\usetheme{metropolis}

\graphicspath{{../../img}}

\begin{document}

\AtBeginSection[]{
  \begin{frame}{Exemples d'applications de la programmation parallèle}
  \tableofcontents[currentsection]
  \end{frame} 
}

\title[]{Exemples d'applications de la programmation parallèle} 
\subtitle{Module 10 \\
INF8601 Systèmes informatiques parallèles}
\author{Sébastien Darche, Michel Dagenais, \\\texttt{< sebastien.darche} \textit{at} \texttt{polymtl.ca >}} 
\date{\today} 
\institute{Polytechnique Montréal}

\begin{frame}[plain]
  \titlepage
\end{frame}

\begin{frame}{Sommaire}
  \tableofcontents
\end{frame}

\section{}

\begin{frame}{Exemples d'applications}

  \begin{itemize}
    \item Décomposition à fine granularité avec organisation MapReduce

    \item Préfixe parallèle

    \item Calcul matriciel

    \item Tri parallèle
  \end{itemize}
\end{frame}

\section{Map Reduce}

\begin{frame}{Map Reduce}

  \begin{itemize}
    \item Le calcul parallèle suivi d'une réduction est un patron commun, bien supporté par OpenMP, MPI...

    \item Cadriciels génériques pour le calcul parallèle basés sur ce paradigme: Google Map Reduce, Hadoop, Apache Spark...

    \item Idéalement, chaque serveur de calcul possède déjà les données à traiter grâce à un système de fichiers distribué comme HFS (Hadoop File System). 

    \item Apache Spark évite de tout écrire sur disque, par exemple pour le traitement en ligne, et peut dans certains cas être beaucoup plus rapide que Hadoop.
  \end{itemize}
\end{frame}

\begin{frame}{Les phases de Map Reduce}

  \begin{itemize}
    \item L'entrée est décomposée en morceaux selon des critères implicites (e.g. ordre d'arrivée) ou explicites (e.g. localisation spatiale, intervalles de numéro d'identification, URL).

    \item La même fonction (programme) est appliquée (map) à toutes les entrées et produit une sortie clé-valeur.

    \item Les sorties peuvent être envoyées vers les noeuds pour la réduction selon les clés (shuffle). 

    \item La réduction (reduce) combine les sorties pour produire le résultat final.
  \end{itemize}
\end{frame}

\begin{frame}[fragile]{Exemple: compter les mots avec Java/Hadoop}

  \scriptsize
  \begin{minted}{java}
// Tutorial from https://hadoop.apache.org/

public class WordCount {

  public static class TokenizerMapper
       extends Mapper<Object, Text, Text, IntWritable>{

    private final static IntWritable one = new IntWritable(1);
    private Text word = new Text();

    public void map(Object key, Text value, Context context
                    ) throws IOException, InterruptedException {
      StringTokenizer itr = new StringTokenizer(value.toString());
      while (itr.hasMoreTokens()) {
        word.set(itr.nextToken());
        context.write(word, one);
      }
    }
  }
  \end{minted}
\end{frame}

\begin{frame}[fragile]{Exemple: compter les mots avec Java/Hadoop}

  \small
  \begin{minted}{java}
  public static class IntSumReducer
       extends Reducer<Text,IntWritable,Text,IntWritable> {
    private IntWritable result = new IntWritable();

    public void reduce(Text key, Iterable<IntWritable> values,
                       Context context
                       ) throws IOException, InterruptedException {
      int sum = 0;
      for (IntWritable val : values) {
        sum += val.get();
      }
      result.set(sum);
      context.write(key, result);
    }
  }
  \end{minted}
\end{frame}

\begin{frame}{Tolerance aux pannes}

  \begin{itemize}
    \item Chaque morceau de fichier est présent en plusieurs (e.g. 3) copies.

    \item Plusieurs entrées peuvent être envoyées à chaque noeud dans la grappe de calcul parallèle.

    \item Si un noeud tombe en panne, les entrées peuvent être réparties sur d'autres noeuds.

    \item Le délai encouru en cas de panne demeure faible.
  \end{itemize}
\end{frame}

\section{Préfixe parallèle}

\begin{frame}{Le problème}

  \begin{itemize}
    \item Une opération associative mais non commutative est appliquée à une suite d'éléments, et produit en sortie une autre suite, dont chaque élément dépend de tous les éléments précédents ou égal en entrée (e.g. $a_i = \sum_{j=1}^i a_j$).

    \item Problème très simple à programmer dans une boucle mais difficile à paralléliser.

    \item Quelques stratégies de parallélisation existent comme procéder en deux passes ou faire un arbre de calcul.

    \item MPI offre la fonction MPI\_Scan avec un choix d'opérations à appliquer comme addition, maximum, minimum...

    \item TBB offre parallel\_scan.
  \end{itemize}
\end{frame}

\begin{frame}[fragile]{Algorithme séquentiel}

  \scriptsize
  \begin{minted}{c}
void seqprfsum(int *u, int m) { 
  int s = u[0];

  for(int i = 1; i < m; i++) { 
    u[i] += s; 
    s = u[i]; 
  } 
}

  \end{minted}
\end{frame}

\begin{frame}[fragile]{Préfixe parallèle en 2 passes}

  \begin{itemize}
    \item Les données sont divisées en plusieurs morceaux à traiter en parallèle (un morceau par thread ou par noeud).

    \item Lors d'une première passe, l'opération est effectuée localement, sans tenir compte du préfixe venant des morceaux précédents.

    \item Ensuite, une réduction est effectuée pour calculer le préfixe pour chaque morceau.

    \item Lors de la seconde passe, l'opération est effectuée à nouveau sur chaque morceau mais cette fois-ci en tenant compte du préfixe venant du morceau précédent.
  \end{itemize}
\end{frame}

\begin{frame}[fragile]{Exemple de calcul en 2 passes avec OpenMP}

  \scriptsize
  \begin{minted}{c}
void parprfsum (int *x, int n, int *z)
{ 
  #pragma omp parallel
  { 
    int id = omp_get_thread_num ();
    int nth = omp_get_num_threads(), chunksize = n / nth;
    int start = id * chunksize;
    seqprfsum(&x[start], chunksize);

    #pragma omp barrier
    #pragma omp single
    { 
      for(int i = 0; i < nth - 1; i++) 
        z[i] = x[(i+1) * chunksize - 1];
      seqprfsum(z ,nth -1);
    }
    if(id > 0 ) 
      for(int j = start; j < start + chunksize; j ++)
        x[j] += z[id - 1];
  }
}
  \end{minted}
\end{frame}

\begin{frame}[fragile]{Exemple de décodage RLE}

  \scriptsize
  \begin{minted}{c}
void uncomprle(int *x, int nx, int *tmp, int *y, int *ny)
{ int nx2 = nx / 2;
  int z[MAXTHREADS];
  tmp[0] = 0;

  for(int i = 0; i < nx2; i++) 
    tmp[i + 1] = x[2 * i];
  
  parprfsum(tmp + 1, nx2 + 1, z);

  #pragma omp parallel
  {
    #pragma omp for
    for(int j = 0; j < nx2; j++) {
      int start = tmp[j];
      int val = x[2 * j +1];
      int nrun = x[2 * j];
      for(int k = 0; k < nrun; k++) 
        y[start + k] = val;
    }
  }
  *ny = tmp[nx2];
}
  \end{minted}
\end{frame}

\begin{frame}[fragile]{Calcul en arbre}

  \begin{itemize}
    \item Dans la barrière en \textit{Papillon}, chaque processus envoie à un autre à chaque $\log_2(n)$ étape (i+1 \% n, i+2 \% n, i+4 \% n...). 

    \item Somme en \textit{papillon}, chaque noeud ajoutant un élément de la séquence (plutôt que 1 pour une barrière).

    \item Si on veut la somme courante ($\sum_{j=1}^i a_i$), plutôt que la somme totale, il suffit de sauter des étapes. On peut vérifier par substitution que $x_6 = (x_0 + (x_1+x_2)) + ((x_3+x_4) + (x_5+x_6))$
  \end{itemize}
  \begin{columns}[t]
    \begin{column}{0.33\textwidth}
$x_1 \leftarrow x_0 + x_1$ \\
$x_2 \leftarrow x_1 + x_2$ \\
$x_3 \leftarrow x_2 + x_3$ \\
$x_4 \leftarrow x_3 + x_4$ \\
$x_5 \leftarrow x_4 + x_5$ \\
$x_6 \leftarrow x_5 + x_6$ \\
$x_7 \leftarrow x_6 + x_7$ \\
    \end{column}
    \begin{column}{0.33\textwidth}
$x_2 \leftarrow x_0 + x_2$ \\
$x_3 \leftarrow x_1 + x_3$ \\
$x_4 \leftarrow x_2 + x_4$ \\
$x_5 \leftarrow x_3 + x_5$ \\
$x_6 \leftarrow x_4 + x_6$ \\
$x_7 \leftarrow x_5 + x_7$ \\
    \end{column}
    \begin{column}{0.33\textwidth}
$x_4 \leftarrow x_0 + x_4$ \\
$x_5 \leftarrow x_1 + x_5$ \\
$x_6 \leftarrow x_2 + x_6$ \\
$x_7 \leftarrow x_3 + x_7$ \\
    \end{column}
  \end{columns}
\end{frame}

\section{Calcul matriciel en parallèle}

\begin{frame}{Le calcul matriciel}

  \begin{itemize}
    \item Plusieurs applications avec beaucoup de données sont exprimées sous forme de calcul matriciel.

    \item Simulation de phénomènes physiques en 2D ou 3D.

    \item Traitement d'image et de vidéo.

    \item Calcul sur des graphes, par exemple représentant des réseaux routiers ou électriques.

    \item Traitement de statistiques.

    \item Apprentissage machine.
  \end{itemize}
\end{frame}

\begin{frame}{Opérations courantes}

  \begin{itemize}
    \item Multiplication de matrices.

    \item Inversion de matrices.

    \item Calcul de déterminant et de valeurs propres.

    \item Puissance de matrices ($M^n$).

    \item Les matrices éparses sont un cas particulier pour lequel plusieurs raccourcis sont disponibles pour en réduire la taille et l'effort de calcul.
  \end{itemize}
\end{frame}

\begin{frame}{La multiplication de matrices}

  \begin{itemize}
    \item Chaque élément du produit de deux matrices carrées de taille n ($n^2$ éléments) demande $n$ multiplications et additions ($\mathcal{O}(n^3)$).

    \item Le calcul d'un élément $c_{ij}$ de la matrice $C = AB$ demande la rangée $i$ de $A$ et la colonne $j$ de $B$.

    \item Si on a $p$ noeuds disponibles, formant une matrice carrée de $\sqrt{p}$, il faut décomposer la matrice de taille $n$ en sous-matrices de taille $m = n/\sqrt{p}$ (ou en groupes de $n/p$ rangées).

    \item Chaque noeud s'occupe alors d'une sous-matrice de taille $m$ et accède les $m$ rangées de $A$ et colonnes de $B$ correspondantes.

    \item On peut envoyer les matrices complètes $A$ et $B$ à chaque noeud mais il est mieux de faire circuler des morceaux de $A$ dans les rangées et de $B$ dans les colonnes de la grappe de calcul.  
  \end{itemize}
\end{frame}

\begin{frame}[fragile]{Exemple multiplication de matrice par rangées}

% https://github.com/suraj-deshmukh/cannon-algorithm-in-c-using-mpi
  \scriptsize
  \begin{minted}{c}
int rank, size, row, column, count, i, j, k;
char ch;
float *A, *B, *C, a, b, c, n;
...
MPI_Init(NULL, NULL);
MPI_Comm_size(MPI_COMM_WORLD, &size);
MPI_Comm_rank(MPI_COMM_WORLD, &rank);

if(rank==0) { /* Lire A et B carrées de taille size = row*row */ }

MPI_Bcast(&row, 1, MPI_INT, 0, MPI_COMM_WORLD);
int periods[] = {1, 1};
int dims[] = {row, row};
int coords[2]; /* 2 dimensions */
int right = 0, left = 0, down = 0, up = 0; /* voisins */
MPI_Comm cart_comm;
MPI_Cart_create(MPI_COMM_WORLD, 2, dims, periods, 1, &cart_comm );
  \end{minted}
\end{frame}

\begin{frame}[fragile]{Exemple multiplication de matrice par rangées}

  \scriptsize
  \begin{minted}{c}
MPI_Scatter(A, 1, MPI_FLOAT, &a, 1, MPI_FLOAT, 0, cart_comm);
MPI_Scatter(B, 1, MPI_FLOAT, &b, 1, MPI_FLOAT, 0, cart_comm);
MPI_Comm_rank(cart_comm, &rank);
MPI_Cart_coords(cart_comm, rank, 2,coords);
MPI_Cart_shift(cart_comm, 1, coords[0], &left, &right);
MPI_Cart_shift(cart_comm, 0, coords[1], &up, &down);
MPI_Sendrecv_replace(&a,1,MPI_FLOAT,left,11,right,11,cart_comm,MPI_STATUS_IGNORE);
MPI_Sendrecv_replace(&b,1,MPI_FLOAT,up,11,down,11,cart_comm,MPI_STATUS_IGNORE);
c = c + a * b;
for(i = 1; i < row; i++) {
  MPI_Cart_shift(cart_comm, 1, 1, &left, &right);
  MPI_Cart_shift(cart_comm, 0, 1, &up, &down);
  MPI_Sendrecv_replace(&a,1,MPI_FLOAT,left,11,right,11,cart_comm,MPI_STATUS_IGNORE);
  MPI_Sendrecv_replace(&b,1,MPI_FLOAT,up,11,down,11,cart_comm,MPI_STATUS_IGNORE);
  c = c + a * b;
}
MPI_Gather(&c, 1, MPI_FLOAT, C, 1, MPI_FLOAT, 0, cart_comm);
if(rank == 0) { /* Imprimer le résultat C */ }
MPI_Finalize();
  \end{minted}
\end{frame}

\begin{frame}[fragile]{Autres opérations matricielles}

  \begin{itemize}
    \item Mettre une matrice à la puissance $n$ peut se faire en combinant avantageusement des multiplications, par exemple $A^8 = ((A^2)^2)^2$.

    \item Inverser une matrice ou solutionner un système d'équations linéaires $Ax = b$, $x = A^{-1}b$
    \begin{itemize}
      \item Formule de Laplace: $A^{-1}={\frac {1}{\det A}}{}^{\operatorname {t} }{\operatorname {com} A}$
      \item Elimination de Gauss-Jordan ou Décomposition LU
      \item Algorithme itératif de Jacobi
    \end{itemize}
    \item Les valeurs propres sont difficiles à calculer analytiquement. Elle peuvent se calculer de manière itérative par multiplication, comme avec la méthode de Jacobi $b_{k+1} = \frac{Ab_k}{\|Ab_k\|}$.
  \end{itemize}
\end{frame}

\begin{frame}[fragile]{Inversion par la méthode itérative de Jacobi}

  \begin{itemize}
    \item $b = Ax$ (b vecteur de valeurs, A matrice de coefficients, x vecteur de variables)
    \item $x = A^{-1}b$.
    \item $D$ une matrice avec la diagonale de $A$ et $O$ une matrice en enlevant la diagonale de A
    \item $x_{k+1} = D^{-1}(b - Ox_k)$. 
  \end{itemize}
\end{frame}

\section{Le tri en parallèle}

\begin{frame}{Les algorithmes classiques de tri}

  \begin{itemize}
    \item Le tri rapide (quicksort) et le tri fusion (merge sort) sont faciles à paralléliser. La phase initiale du premier (séparer les éléments par le pivot) et la phase finale du second (fusionner à la racine) constituent un goulot d'étranglement. 

    \item Le tri par comparaison, comme le tri en bulle, normalement moins efficace en série, n'est pas si mal en parallèle car plus équilibré.

    \item Le tri par paquets demeure très efficace en parallèle, si on choisit de bons seuils pour décomposer en paquets. Une phase d'échantillonage permet d'estimer la distribution et choisir de bons seuils.
  \end{itemize}
\end{frame}

\section{Conclusion}

\begin{frame}{Conclusion}

  \begin{itemize}
    \item Il est intéressant de connaître un certain nombre de stratégies classiques de parallélisation pour des problèmes courants.

    \item Cependant, chaque plate-forme (OpenMP, OpenCL, MPI) a des particularités différentes qui affectent l'efficacité des algorithmes.

    \item Les particularités des jeux de données, pour les très grands ensembles, peuvent aussi affecter le choix de l'algorithme (e.g. matrices éparses).

    \item Il existe une littérature scientifique étendue sur les meilleurs algorithmes pour un très large éventail de problèmes.
  \end{itemize}
\end{frame}

\end{document}
